\chapter{Experiments}\label{chap:experiments}

\section{Overview}\label{sec:experiments-overview}

This chapter reports what both phases actually found. Phase~1 (Section~\ref{sec:phase1}) covers the regression work on the pre-war corpus: the headline bias experiment, the contemporaneous and lag OLS regressions that give the first evidence for RQ1 and RQ2, and the VAR. Phase~2 (Section~\ref{sec:phase2}) covers the deep-learning work on the expanded corpus: the LLM extraction pipeline, the two Temporal Fusion Transformers, and the ablation that shaped the v2 design. The chapter ends by comparing the two phases as methodologies (Section~\ref{sec:phase-comparison}), setting out what Phase~2 changed and why putting the two phases' numbers side by side would be a mistake.

\section{Phase 1: initial pipeline and baseline modeling}\label{sec:phase1}

Phase~1 builds the empirical baseline, starting from the existing evidence that news sentiment predicts trading activity and volume rather than just prices \citep{tetlock2007}. It runs on the FinBERT sentiment of Section~\ref{sec:finbert} and the regex body-validity filter of Section~\ref{sec:regex-filter}, over WTI-relevant news published between March 2024 and February 2026. Three techniques get applied in turn: linear regression (contemporaneous and lagged), vector autoregression, and a headline-bias comparison that is what pushes us toward using full article bodies later. By the end of Phase~1 we have the canonical statistical evidence for RQ1 (lag structure) and RQ2 (bearish-versus-bullish asymmetry), and the limits of a FinBERT-plus-classical-models pipeline are clear enough to justify the Phase~2 changes.

\subsection{Phase 1 data and feature setup}\label{sec:phase1-data}

The Phase~1 modeling dataset is \textbf{13{,}690 articles}: 7{,}756 with substantive bodies (regex filter \texttt{body\_valid=1}, Section~\ref{sec:regex-filter}) feeding title-plus-body FinBERT sentiment, and 5{,}934 where the body fetch failed and the title was kept as a fallback. Behind that are 51{,}948 raw GDELT records across eight oil-market queries, cut to 16{,}326 unique English articles after deduplication (Table~\ref{tab:phase1-data}). FinBERT \citep{araci2019} scores each article (Section~\ref{sec:finbert}); the regressions use the title-plus-body class probabilities as the continuous sentiment feature, with the discrete $\{+1, 0, -1\}$ label kept as a baseline. The sentiment splits \textbf{44.9\% bearish, 30.4\% bullish, 24.7\% neutral}. The bearish skew is a property of headline tone, documented in Section~\ref{sec:headline-bias}, and should not be read as a claim about the population.

Articles get aligned to trading hours by the ceiling rule (Section~\ref{sec:ceiling}), which puts 82.7\% within two hours of a trading hour (contemporaneous) and forward-assigns the rest. The market backbone is 11{,}219 hourly WTI records carrying the liquidity targets (\texttt{log\_volume}, \texttt{price\_range}, \texttt{amihud} \citep{amihud2002}) from Section~\ref{sec:data-sources}. DXY, VIX, and the EIA covariates only show up in Phase~2. Table~\ref{tab:phase1-data} summarises how the dataset was built.

\begin{table}[h]
  \centering
  \caption{Phase 1 dataset construction.}
  \label{tab:phase1-data}
  \begin{tabular}{@{}p{9cm}l@{}}
    \toprule
    Stage & Count \\
    \midrule
    Raw GDELT articles & 51{,}948 \\
    After deduplication + English filter & 16{,}326 \\
    With successful body retrieval & $\sim$13{,}000 (80\% of 16{,}326) \\
    With \texttt{body\_valid=1} (regex-passed body) & 7{,}756 \\
    Title-only fallback (body failed but title retained) & 5{,}934 \\
    \textbf{Phase 1 modeling dataset} & \textbf{13{,}690} \\
    Aligned market hours (raw, before VAR cleaning) & 11{,}219 \\
    Articles aligned within 2h gap (contemporaneous) & 11{,}316 \\
    Articles aligned with gap $\geq$ 2h (forward-assigned) & 2{,}374 \\
    Article date range & 2024-03-11 to 2026-03-01 \\
    \bottomrule
  \end{tabular}
\end{table}

Everything in the three Phase~1 experiments runs on this dataset.

\subsection{Headline bias experiment}\label{sec:headline-bias}

\textbf{Motivation.} The Phase~1 pipeline mixes two FinBERT input regimes in one 13{,}690-article dataset: title-only for the fallback articles, and title-plus-body for the rest. Before treating all of that as a single signal, we should check whether the two regimes actually agree. The question is also interesting on its own. Headlines exist to get attention, and a body can qualify or even reverse whatever the title implies.

\textbf{Method and categorical result.} We took the 7{,}755 articles that have both a title and a substantive body and ran FinBERT twice on each, once on the title alone and once on title-plus-body. Because the model is the same both times, any difference comes from the input content (Section~\ref{sec:finbert}). On the argmax labels the two \textbf{disagree on 41.6\% of articles} ($\chi^2 = 2050$, $p < 0.001$), and the disagreement has structure to it: positive titles keep a positive body only 59.2\% of the time, with 40.8\% moving toward neutral or negative; neutral titles resolve negative more often than positive (39.6\% vs 30.0\%); negative titles stay negative in 78.2\% of cases.

\textbf{Continuous result.} The label measure has a blind spot. It ignores how far sentiment moves, and it cannot see any shift that stays inside one label. So we also define a signed score $P(\text{positive}) - P(\text{negative})$ per input and a divergence magnitude $\lvert \text{signed\_full} - \text{signed\_title}\rvert$. Label flips carry a mean magnitude of \textbf{0.96}, which is close to a full reversal. Same-label articles still move by \textbf{0.17} on average, with 31.6\% moving more than 0.20. The mean signed shift from title to body is \textbf{$-0.09$}, and 57.6\% of articles get more bearish once the body is read.

\textbf{Worked example.} These quantities are easier to follow on real articles. The signed score is $P(\text{positive}) - P(\text{negative})$, computed once on the title and once on title plus body. The magnitude is the absolute distance between them, so how far sentiment moved, and the signed shift is $\text{signed\_full} - \text{signed\_title}$, so which way it moved.

\begin{table}[h]
  \centering
  \small
  \caption{Worked examples of title-versus-body sentiment divergence.}
  \label{tab:headline-example}
  \begin{tabular}{@{}p{4.5cm}rrrrp{2.3cm}@{}}
    \toprule
    Article (title) & signed\_title & signed\_full & magnitude & signed shift & Label transition \\
    \midrule
    ``Oil price soars above \$80'' & $+0.49$ & $-0.63$ & 1.12 & $-1.12$ & positive $\to$ negative (flip) \\
    ``U.S. Oil Is Stealing Market Share from OPEC+'' & $-0.32$ & $-0.77$ & 0.45 & $-0.45$ & negative $\to$ negative \\
    ``Oil prices continue to drop in world markets'' & $-0.96$ & $-0.96$ & 0.00 & 0.00 & negative $\to$ negative \\
    \bottomrule
  \end{tabular}
\end{table}

The first article flips. FinBERT reads the headline as bullish (``soars above \$80'') and then turns bearish once it sees the body, a near-full reversal at magnitude 1.12. The second keeps its negative label but still moves 0.45 further bearish, which is exactly the within-bucket movement the 41.6\% flip rate cannot count. The third does not move at all, title and body agree. In all three the body is at least as bearish as the title, which is the one-directional pattern that produces the $-0.09$ mean signed shift.

\textbf{Divergence and label-flip totals.} Table~\ref{tab:headline-totals} collects the aggregate quantities across the 7{,}755 title-plus-body articles.

\begin{table}[h]
  \centering
  \caption{Title-versus-body divergence totals (7{,}755 articles).}
  \label{tab:headline-totals}
  \begin{tabular}{@{}lr@{}}
    \toprule
    Quantity & Value \\
    \midrule
    Label agreement (title and body same label) & 4{,}527 (58.4\%) \\
    Label flip (disagreement) & 3{,}228 (41.6\%) \\
    Flip rate, positive titles & 40.8\% \\
    Flip rate, neutral titles & 69.5\% \\
    Flip rate, negative titles & 21.8\% \\
    Mean divergence magnitude, all articles & 0.50 \\
    Mean divergence magnitude, label flips & 0.96 \\
    Mean divergence magnitude, same-label & 0.17 \\
    Same-label articles with magnitude above 0.20 & 31.6\% \\
    Mean signed shift, title to body & $-0.09$ \\
    Articles more bearish once the body is read & 57.6\% \\
    \bottomrule
  \end{tabular}
\end{table}

\textbf{Implications.} Two things carry forward. The first is that the input regime is not a detail: a pipeline that swaps title-only and title-plus-body freely is averaging over a 41.6\% disagreement rate. The second cuts against the usual intuition. People expect attention-grabbing headlines to overstate negativity, but here \textbf{titles lean more bullish than the articles under them}, and reading the body pulls FinBERT's sentiment more negative on average. Only the signed continuous score establishes that direction; the labels alone cannot. So the 5{,}934 title-only fallback articles carry a sentiment that would have moved in roughly 40\% of cases if a body had been available. Since that shift goes one way across all lags, it weakens the OLS effect reported below rather than inflating it. Beyond Phase~1, this is the first concrete evidence that the choice of news representation matters, and it is what motivates the Phase~2 move to structured extraction from full bodies (Section~\ref{sec:schema}).

\subsection{Lag OLS: RQ1 and RQ2}\label{sec:lag-ols}

The lag OLS regressions are the canonical Phase~1 evidence for RQ1 (the lag structure of news impact on liquidity) and RQ2 (bearish-versus-bullish asymmetry). They are also the result that the Phase~2 methodology is ultimately trying to corroborate.

\textbf{Specification.} At horizon $k$,
\begin{equation}
  \text{log\_volume}_{t+k} = \beta_0 + \beta_1\, P(\text{negative})_t + \beta_2\, P(\text{positive})_t + \varepsilon_t ,
\end{equation}
where $P(\text{negative})_t$ and $P(\text{positive})_t$ are the title-plus-body FinBERT class probabilities for the article published at hour $t$, and the neutral probability is the omitted reference (Section~\ref{sec:ols-spec}). So $\beta_1$ and $\beta_2$ give the \texttt{log\_volume} response to a maximally confident bearish or bullish article against a neutral one. Each article is one observation, and we estimate a separate regression at each $k \in \{0, 1, 2, 3, 4, 6, 8, 12\}$, with $k = 0$ the contemporaneous case. Sample size runs from 13{,}690 rows at $k = 0$ down to roughly 9{,}900 at the longest lag. The continuous encoding is primary. The discrete-dummy baseline is kept for comparison, and at $k = 0$ it gives smaller coefficients ($\beta_{\text{bearish}} = 0.133$, $\beta_{\text{bullish}} = 0.103$) than the continuous 0.186 / 0.166, which confirms that weighting each article by FinBERT's confidence sharpens the signal.

\begin{table}[h]
  \centering
  \caption{Lag OLS coefficients across horizons.}
  \label{tab:lag-ols}
  \begin{tabular}{@{}lrrrrrr@{}}
    \toprule
    Lag $k$ (h) & $\beta_{P(\text{neg})}$ & $p$ & $\beta_{P(\text{pos})}$ & $p$ & $R^2$ & $n$ \\
    \midrule
    0 (contemp.) & 0.186 & 0.000 & 0.166 & 0.000 & 0.0015 & 13{,}690 \\
    1 & 0.262 & 0.000 & 0.200 & 0.000 & 0.0028 & 11{,}598 \\
    2 & 0.105 & 0.068 & 0.114 & 0.084 & 0.0003 & 11{,}512 \\
    3 & 0.117 & 0.047 & 0.170 & 0.013 & 0.0006 & 11{,}360 \\
    4 & 0.267 & 0.000 & 0.237 & 0.001 & 0.0017 & 11{,}065 \\
    6 & 0.342 & 0.000 & 0.291 & 0.000 & 0.0027 & 10{,}679 \\
    8 & 0.056 & 0.409 & 0.101 & 0.192 & 0.0002 & 10{,}440 \\
    12 & 0.001 & 0.984 & $-0.171$ & 0.029 & 0.0008 & 9{,}917 \\
    \bottomrule
  \end{tabular}
\end{table}

\textbf{Result.} Table~\ref{tab:lag-ols} has the regression at every horizon, one row per lag $k$ in hours, with $k = 0$ the contemporaneous case. Reading the columns: $\beta_{P(\text{neg})}$ and $\beta_{P(\text{pos})}$ are the fitted coefficients on the bearish and bullish class probabilities, meaning the expected change in \texttt{log\_volume} for a maximally confident bearish or bullish article against a neutral one. The two $p$ columns are their two-sided p-values, with the usual significance threshold at $p < 0.05$. $R^2$ is the share of \texttt{log\_volume} variance the two-feature regression explains at that lag. And $n$ is how many articles enter that regression, which shrinks as $k$ grows because targets at hour $t + k$ past the coverage window get dropped.

\begin{figure}[h]
  \centering
  \includegraphics[width=\textwidth]{lag_coefficients}
  \caption{Lag OLS coefficients and p-values across horizons. Left panel: the bearish $P(\text{negative})$ and bullish $P(\text{positive})$ coefficients at lags 1, 2, 3, 4, 6, 8, and 12 hours. Stars mark lags significant at the conventional $p < 0.05$ threshold. Right panel: the p-values for the same coefficients, with the $p = 0.05$ reference line. The peak at lag $+6$h is the canonical finding for RQ1; the ordering of bearish above bullish at the dominant lags is the canonical finding for RQ2.}
  \label{fig:lag-coefficients}
\end{figure}

\textbf{RQ1 (lag structure).} The impact is detectable straight away at publication ($k = 0$), builds over the first few hours, peaks sharply at \textbf{$+6$h} ($\beta_{P(\text{neg})} = 0.342$, $\beta_{P(\text{pos})} = 0.291$, both $p < 0.001$), and has decayed to insignificance by $+8$h. To put the peak in concrete terms, the $+6$h bearish coefficient means a fully confident bearish article implies roughly 41\% more volume than a neutral hour, since $\exp(0.342) - 1 \approx 0.41$. The shape is not a smooth exponential decay. There are secondary maxima at $+1$h and $+4$h, which points to discrete information-incorporation events at particular horizons rather than a continuous fade.

\textbf{RQ2 (asymmetry).} Bearish articles get a stronger response than bullish ones at the lags that matter (0, 1, 4, and the $+6$h peak, where bearish beats bullish by about 18\%). Lag 3 is the exception, with both coefficients small and bullish marginally ahead, but it is a weak lag and it does not overturn the pattern at the horizons carrying the effect. Since the two probabilities are estimated as independent regressors (Section~\ref{sec:ols-spec}), the bearish-over-bullish ordering is something the data produced, not something the model was set up to find.

\textbf{On the small $R^2$ and robustness.} $R^2$ never gets above 0.003. That is what we expected. Hourly volume has many drivers and sentiment is one of them, so $R^2$ here measures sentiment's share of \emph{total} volume variance, not of \emph{news-driven} volume. The finding is the structured coefficient pattern, not the explained variance. Two robustness notes. The regression is cross-sectional, one article per row with no strong serial correlation, so we do not apply HAC corrections (Section~\ref{sec:ols-spec}). And the ceiling alignment guarantees every article precedes its matched hour, which if anything understates the immediate reaction, so the bias runs conservative for the peak-lag finding. Phase~2 revisits both RQ1 and RQ2 independently through the TFT v2 attention pattern (Section~\ref{sec:tft-v2}).

\subsection{Phase 1 summary}\label{sec:var}

Phase~1 produces the canonical statistical evidence for both research questions, and it exposes the limitations Phase~2 goes on to address.

\textbf{Evidence.} RQ1 gets its answer from the lag OLS peak at \textbf{$+6$h}, and RQ2 from the \textbf{bearish $>$ bullish} ordering at the dominant lags (Section~\ref{sec:lag-ols}). Both are the primary Phase~1 evidence, and the TFT v2 attention pattern revisits both independently in Phase~2. The headline-bias experiment (Section~\ref{sec:headline-bias}) adds a separate finding: title-only and title-plus-body FinBERT disagree on 41.6\% of articles, with titles leaning more bullish than the bodies they head, which is what motivates extracting features from full bodies in Phase~2.

\textbf{VAR exploration (abandoned).} We fitted an exploratory VAR and then abandoned it. More than half the hours carry no contemporaneous news, so the sentiment series is mostly zeros and the impulse responses came back insignificant. The failure is still informative: it tells us the data are sparse and event-driven, and that property is exactly what motivates the deep-learning model in Phase~2.

\textbf{Limitations that drive Phase 2.} Five limitations of the Phase~1 pipeline each map onto a Phase~2 change (Table~\ref{tab:phase1-limitations}).

\begin{table}[h]
  \centering
  \caption{Phase 1 limitations and the Phase 2 responses.}
  \label{tab:phase1-limitations}
  \begin{tabular}{@{}p{7cm}p{6.5cm}@{}}
    \toprule
    Phase 1 limitation & Phase 2 response \\
    \midrule
    FinBERT reduces each article to one sentiment axis (no magnitude, event type, certainty, horizon, entities) & Structured LLM extraction with a richer schema (Section~\ref{sec:llm-extraction}) \\
    The regex \texttt{body\_valid} filter is fast but not semantically discriminative (false positives and negatives) & LLM-judged \texttt{usable} flag (Section~\ref{sec:filter-comparison}) \\
    The sentiment signal is sparse ($\sim$50\% of hours have no news), defeating models that need joint observation & A TFT built for sparse, event-driven inputs (Section~\ref{sec:tft-v1}) \\
    The corpus is pre-war only (ends February 2026) & Corpus extended through May 2026 (Section~\ref{sec:phase2-data}) \\
    The OLS has no macro controls & DXY and VIX added as exogenous covariates (Section~\ref{sec:tft-v1}) \\
    \bottomrule
  \end{tabular}
\end{table}

Every Phase~2 step answers a concrete limitation the previous step exposed, so the progression can be traced rather than just claimed.

\section{Phase 2: refined pipeline and deep-learning modeling}\label{sec:phase2}

Phase~2 takes on the Phase~1 limitations (Section~\ref{sec:var}) with a coordinated redesign: FinBERT out and structured LLM extraction in, DXY and VIX added as macro covariates, the regex filter replaced by an LLM-judged \texttt{usable} flag, a Temporal Fusion Transformer as the primary model, and the corpus extended through May 2026. It happened in a particular order, and the write-up follows that order. An initial LLM schema (v1) supported the first TFT training (TFT v1). An inter-model calibration then exposed a weakness in the composite sentiment score, which motivated a revised schema (v2) with three orthogonal economic channels and the \texttt{usable} flag. The full extraction then ran on the expanded corpus and supported the final TFT v2. So the subsections go: data setup (Section~\ref{sec:phase2-data}), LLM extraction (Section~\ref{sec:llm-extraction}), filter migration (Section~\ref{sec:filter-comparison}), first calibration (Section~\ref{sec:calibration}), TFT v1 (Section~\ref{sec:tft-v1}), the schema response (Section~\ref{sec:channel-response}), TFT v2 (Section~\ref{sec:tft-v2}), and the comparison of the two phases (Section~\ref{sec:phase-comparison}).

\subsection{Phase 2 data and feature setup}\label{sec:phase2-data}

The Phase~2 corpus extends Phase~1 two ways: the GDELT scrape runs through May 2026 to cover the post-war period, and five new queries go after geopolitical and demand-side themes (Section~\ref{sec:data-sources}). The full corpus is \textbf{22{,}795 article records} from January 2024 to May 2026. Of those, \textbf{19{,}619 have substantive body content} (non-null, not a scraping error) and go to LLM extraction, while 3{,}176 came back with null bodies from a paywall, a fetch failure, or JavaScript rendering.

\textbf{Macro covariates.} We added hourly DXY (\texttt{DX-Y.NYB}) and VIX (\texttt{\^{}VIX}) series as exogenous controls (Section~\ref{sec:input-typing}). The reason is empirical: the Phase~1 OLS explained under 0.3\% of volume variance (Section~\ref{sec:lag-ols}), so letting the model soak up broad market-driven variation should leave a cleaner residual for news to explain. We also moved the pipeline off CSV intermediates and onto the SQLite database (\texttt{wti\_thesis.db}), which makes the alignment a deterministic join.

\textbf{Aligned dataset.} Re-running the temporal alignment (Section~\ref{sec:temporal-alignment}) over the expanded corpus gives \textbf{22{,}795 article-aligned rows} (15{,}290 contemporaneous, 7{,}505 forward-assigned) on the 11{,}232-hour market grid from 13 May 2024 to 13 May 2026. Table~\ref{tab:phase-comparison-data} puts the two phases side by side.

\begin{table}[h]
  \centering
  \small
  \caption{Phase 1 versus Phase 2 dataset.}
  \label{tab:phase-comparison-data}
  \begin{tabular}{@{}p{3cm}p{4.7cm}p{5.7cm}@{}}
    \toprule
    Stage & Phase 1 & Phase 2 \\
    \midrule
    GDELT queries & 8 & 13 \\
    Article window & March 2024 -- February 2026 & January 2024 -- May 2026 \\
    Raw articles (GDELT scrape) & 51{,}948 & 76{,}345 (cumulative; Phase 2 re-scrape added 24{,}397) \\
    After deduplication + English & 16{,}326 & 22{,}795 \\
    Articles with substantive bodies & 7{,}756 (\texttt{body\_valid=1}) & 19{,}619 (non-null, non-error body) \\
    Title-only fallback & 5{,}934 & 0 (no title-only fallback in Phase 2) \\
    Modeling-ready corpus & 13{,}690 (Phase 1 OLS/VAR) & 10{,}514 (\texttt{usable\_strict}) / 11{,}675 (\texttt{usable}) \\
    Market hours window & March 2024 -- February 2026 & May 2024 -- May 2026 \\
    Hourly market observations & 11{,}219 & 11{,}232 \\
    Macro covariates & none & DXY, VIX \\
    Sentiment representation & FinBERT 3-class probabilities & LLM continuous sentiment + 3 channels + categorical fields \\
    \bottomrule
  \end{tabular}
\end{table}

Everything else in this chapter runs on the Phase~2 dataset.

\subsection{LLM feature extraction}\label{sec:llm-extraction}

Swapping FinBERT for LLM extraction is the central Phase~2 change. It recovers the dimensions a single sentiment axis throws away: magnitude, event type, certainty, time horizon, entities, and the economic channels. Sections~\ref{sec:schema} through~\ref{sec:other-fields} specify the methodology. This subsection is about what happened when we ran it.

\textbf{Extraction setup.} Claude Haiku 4.5 scored each of the 19{,}619 articles with a substantive body, through the tool-use API, which forces the \texttt{extract\_article\_features} schema (Section~\ref{sec:tool-use}) and lets us parse the tool call deterministically. The user message is the title and body, truncated to 1{,}500 characters, and Appendix~\ref{app:prompt} has the full prompt and tool schema. The schema short-circuits unusable articles down to just the \texttt{usable} flag, which cuts output tokens on the roughly 40\% unusable share, and all 19{,}619 articles returned a valid response with zero unrecoverable errors. The deployed schema is v2 (Section~\ref{sec:channel-decomposition}): the composite \texttt{sentiment\_score}, three orthogonal channels (\texttt{supply\_impact}, \texttt{demand\_impact}, \texttt{risk\_premium}), \texttt{magnitude}, \texttt{certainty}, \texttt{event\_type} as an array of one to three categories, \texttt{time\_horizon}, \texttt{entities}, and the required \texttt{usable} flag. The batch filtered on non-null, non-error bodies rather than the Phase~1 regex, which makes the LLM's own \texttt{usable} flag the downstream filter. Section~\ref{sec:filter-comparison} compares it against the regex.

\textbf{Extraction outcome.} The LLM flagged \textbf{11{,}675 articles as \texttt{usable} (59.5\%)} and 7{,}944 as unusable (40.5\%, mostly paywalls, Cloudflare blocks, and keyword-matched off-topic content). That lines up with both the pre-batch estimate and the calibration sample (Section~\ref{sec:calibration}). Table~\ref{tab:extraction-means} gives the four continuous scores across the 11{,}675 usable articles.

\begin{table}[h]
  \centering
  \caption{Continuous feature means over the 11{,}675 usable articles.}
  \label{tab:extraction-means}
  \begin{tabular}{@{}lrr@{}}
    \toprule
    Feature & Mean & Std \\
    \midrule
    \texttt{sentiment\_score} & $+0.025$ & 0.501 \\
    \texttt{supply\_impact} & $-0.071$ & 0.463 \\
    \texttt{demand\_impact} & $-0.075$ & 0.340 \\
    \texttt{risk\_premium} & $+0.156$ & 0.426 \\
    \bottomrule
  \end{tabular}
\end{table}

The near-zero \texttt{sentiment\_score} mean says the corpus is roughly balanced between bearish and bullish. The negative \texttt{supply\_impact} and \texttt{demand\_impact} means say it tilts toward supply-tightening and demand-weakening news. And the positive \texttt{risk\_premium} mean, the largest of the four in magnitude, reflects how elevated geopolitics were over the coverage period. One number worth flagging: 9.9\% of usable articles have all three channels at exactly zero, meaning the LLM called them usable but materially neutral. Section~\ref{sec:calibration} comes back to this, and Section~\ref{sec:filter-comparison} uses it to define the stricter \texttt{usable\_strict} filter.

\textbf{Categorical and entity distributions.} First-rank event types are mostly \texttt{geopolitical} (37.6\%) and \texttt{supply} (28.6\%), then \texttt{macro} (14.5\%) and \texttt{demand} (12.7\%). The most frequent entities are Iran (3{,}197), the United States (3{,}092), Russia, China, OPEC+, Trump, Israel, India, the Strait of Hormuz (832, roughly matching OPEC's count), and OPEC. All of these make sense for the coverage period, and the fact that the Strait of Hormuz shows up as its own entity rather than being folded into Iran is evidence the normalization is separating chokepoints from actors properly. These distributions are the input for the TFT models below.

\subsection{LLM usable flag and filter comparison}\label{sec:filter-comparison}

The Phase~1 regex \texttt{body\_valid} filter (Section~\ref{sec:regex-filter}) and the Phase~2 LLM \texttt{usable} flag (Section~\ref{sec:usable-filter}) are two answers to the same question: which scraped articles carry substantive content? Both labels exist for the 19{,}619 processed articles, so we can just compare them.

\textbf{Result.} Table~\ref{tab:filter-contingency} is the $2 \times 2$ contingency table.

\begin{table}[h]
  \centering
  \caption{Regex versus LLM filter agreement (19{,}619 articles).}
  \label{tab:filter-contingency}
  \begin{tabular}{@{}lrrr@{}}
    \toprule
    & LLM \texttt{usable=0} & LLM \texttt{usable=1} & \textbf{Total} \\
    \midrule
    regex \texttt{body\_valid=0} & 4{,}578 & 1{,}491 & 6{,}069 \\
    regex \texttt{body\_valid=1} & 3{,}366 & 10{,}184 & 13{,}550 \\
    \textbf{Total} & 7{,}944 & 11{,}675 & 19{,}619 \\
    \bottomrule
  \end{tabular}
\end{table}

The regex accepts 69.1\% of articles and the LLM 59.5\%, and they agree on 75.2\% of them (14{,}762 of 19{,}619), which is a \textbf{Cohen's $\kappa$ of 0.47}, chance-corrected, or ``moderate'' on the Landis-Koch convention. The interesting part is the two off-diagonal cells. The regex rejects but the LLM accepts \textbf{1{,}491} articles, the regex's false negatives, which are short but substantive pieces sitting under the 400-character threshold. The regex accepts but the LLM rejects \textbf{3{,}366}, the regex's false positives, which are long keyword-matched articles that are not about WTI. So the regex is too permissive on long non-substantive content, because it cannot judge topic semantically, and too restrictive on short substantive content, because its length cutoff is blunt. The LLM, which can read the body, gets both of these the other way round.

\textbf{Manual audit.} We sampled both off-diagonal cells and hand-labelled them, and the LLM's judgment holds up in the large majority of cases. The regex's false positives are mostly keyword collisions: palm oil, canola, a semiconductor firm with an energy-sounding name, generic equity wraps. Its false negatives are short but on-topic briefs, for instance a 320-character note on OPEC+ holding its production plan. The audit also turned up a residual LLM error mode. A few articles the LLM marks \texttt{usable} are not really WTI-relevant, the clearest being one about Canada rolling back EV tariffs. So the LLM filter beats the regex on both of the regex's failure modes without being a ground-truth oracle itself, a limitation Chapter~\ref{chap:discussion} returns to.

\textbf{The \texttt{usable\_strict} variant.} A distinct group of articles is marked \texttt{usable} and yet comes back with all three channel scores at exactly zero: \textbf{1{,}161} of them, the 9.9\% from Section~\ref{sec:llm-extraction}. We define a stricter filter for these, \texttt{usable\_strict=1}, requiring at least one non-zero channel, which leaves \textbf{10{,}514} articles. It is the canonical filter for TFT v2 training and the reported ablation (Appendix~\ref{app:ablation}), while \texttt{usable=1} is only used for this filter comparison.

Our first hypothesis was that these would turn out to be the LLM's residual false positives, the EV-tariffs kind from the audit above, finally showing themselves through a diagnostic pattern. A manual audit did not confirm it, and the distinction matters for how the filter should be understood. Inspecting the zero-channel articles turned up no off-topic contamination: no EV stories, no scraping artifacts, nothing resembling the keyword collisions that characterise the regex's false positives. They are a mix of two things instead. Some are genuine residual punting, where the LLM tagged the article with a channel-relevant event type (supply, demand, or geopolitical) and then declined to commit to a magnitude on the matching channel, which is the same conservatism the calibration flagged in Section~\ref{sec:channel-response}. The rest are on-topic articles whose channels are correctly zero for what they are: pure price reports, technical analysis, equity-market reactions to oil moves, and macro framing with no specific supply or demand implication. Nothing is wrong with that second group.

So \texttt{usable\_strict} is not a correctness filter and we do not claim it as one. It is a signal filter. The channels are what Phase~2 contributes, and an article carrying no channel signal contributes nothing to them regardless of whether it is a good article, so training on it adds noise to the features the model is meant to learn from. The cost of the choice is that roughly a thousand substantive articles are dropped for a property that is a correct description of them, which is a real trade-off rather than a free improvement, and Appendix~\ref{app:ablation} tests the \texttt{usable} against \texttt{usable\_strict} configurations directly.

\textbf{Filter choice and cross-phase note.} The two filters are not nested. The LLM accepts 1{,}491 articles the regex rejects and rejects 3{,}366 it accepts, so its acceptance set overlaps the regex's without being a subset of it. The LLM filter is the canonical Phase~2 downstream filter. Between this migration and the loss of Phase~1's title-only fallback, the Phase~1 modeling set (13{,}690) and the Phase~2 set (11{,}675 usable, 10{,}514 strict) are not directly comparable, which Section~\ref{sec:phase-comparison} develops.

\subsection{Inter-model calibration of LLM features}\label{sec:calibration}

Since the LLM-extracted features are what all of Phase~2 sits on, we validate them before using them. Following Section~\ref{sec:calibration-methods}, we run an \textbf{inter-model calibration} \citep{gilardi2023, zheng2023}: the same 30 stratified articles get scored by Claude Haiku 4.5 and by a GPT-family reference model with an identical prompt and schema, and we treat cross-family disagreement as a signal of either genuine ambiguity or systematic bias. This subsection reports the first calibration, on Schema v1, and what it told us.

\textbf{Result on Schema v1.} Schema v1 had a single composite \texttt{sentiment\_score} as its directional feature. The two models agreed on the \texttt{usable} flag for 26 of 30 articles (87\%), but agreement on the composite sentiment was bad: a Pearson correlation of only \textbf{0.39}, with \textbf{opposite-signed} scores on \textbf{4 of 13 usable articles (31\%)}. Sign disagreement is the worst kind to have, because it means the two models read the same article as pushing the WTI price in opposite directions.

\textbf{Diagnostic (the Phase 2 pivot).} The sign disagreements were not scattered at random. They piled up on high-magnitude geopolitical events, which are precisely the articles that move the market. The cause is that the composite \texttt{sentiment\_score} is being asked to hold two separate judgments at once: the \textbf{valence of the event} (is this good news or bad news?) and its \textbf{directional price impact} (does it push crude up or down?). For most articles those agree, so nothing breaks. For a supply-threatening escalation they come apart: the event is negative in valence but bullish for price. A model weighting valence scores it negative, a model weighting price scores it positive, and one number cannot hold both. So the two models resolve the ambiguity differently exactly where it matters most. A 0.39 correlation on the primary feature is a structural problem, not bad luck.

\textbf{Response.} The fix is to decompose the judgment rather than fiddle with the prompt. Section~\ref{sec:channel-response} introduces the three orthogonal economic channels (\texttt{supply\_impact}, \texttt{demand\_impact}, \texttt{risk\_premium}) and the re-calibration that validates them. But TFT v1 was trained on the v1 schema before any of that, so its results and limitations come first.

\subsection{Temporal Fusion Transformer (TFT) v1}\label{sec:tft-v1}

TFT v1 is the project's first deep-learning model and the first to use the richer LLM features. It trained on the Phase~1 corpus, with Schema v1 features plus the DXY and VIX macro covariates. The Temporal Fusion Transformer \citep{lim2021} is an attention-based multi-horizon forecaster, and its Variable Selection Network (VSN) and interpretable attention are what make it a good fit here.

\textbf{Training setup.} The architecture (Section~\ref{sec:tft-methods}) uses a 48-hour encoder, a one-hour horizon, \texttt{hidden\_size=32}, \texttt{attention\_head\_size=4}, \texttt{dropout=0.1}, and a quantile loss, for about 113{,}000 parameters. On 10{,}797 hourly rows with an 80/20 temporal split, early stopping fired at epoch 21 with a best validation loss of \textbf{0.204}. The inputs are the Schema v1 LLM features (\texttt{sentiment\_score}, \texttt{magnitude}, integer-encoded \texttt{event\_type}/\texttt{price\_direction}/\texttt{time\_horizon}, \texttt{certainty}, \texttt{n\_articles}), the market context (\texttt{log\_volume}, \texttt{price\_range}, \texttt{log\_return}, \texttt{amihud}, \texttt{dxy}, \texttt{vix}), and calendar covariates.

\textbf{Feature importance.} One feature takes most of the VSN weight (Table~\ref{tab:tftv1-importance}).

\begin{table}[h]
  \centering
  \caption{TFT v1 feature importance.}
  \label{tab:tftv1-importance}
  \begin{tabular}{@{}lrl@{}}
    \toprule
    Feature & Importance & Group \\
    \midrule
    \texttt{sentiment\_score} & 0.53 & News \\
    \texttt{log\_volume} & 0.095 & Market \\
    \texttt{dxy} & 0.075 & Market \\
    \texttt{log\_return} & 0.055 & Market \\
    \texttt{event\_type\_num} & 0.030 & News \\
    \texttt{magnitude} & 0.025 & News \\
    \texttt{vix} & 0.005 & Market \\
    \bottomrule
  \end{tabular}
\end{table}

\begin{figure}[h]
  \centering
  \includegraphics[width=0.85\textwidth]{tft_feature_importance}
  \caption{TFT v1 feature importance, news features (red) versus market features (blue). \texttt{sentiment\_score} dominates at 0.53, more than five times the next feature; the market-context features form a second tier and the rest carry minor weight.}
  \label{fig:tftv1-importance}
\end{figure}

Three things stand out. First, \textbf{\texttt{sentiment\_score} dominates at 0.53, over five times the next feature}, which is a strong result for the LLM approach: the continuous LLM sentiment carries far more signal than any market variable, and far more than a FinBERT label could. Second, \textbf{VIX is negligible at 0.005} while DXY ranks third, so for WTI the dollar index is apparently already absorbing most of the macro-risk signal. Third, the categorical features (\texttt{event\_type}, \texttt{price\_direction}, \texttt{time\_horizon}) were integer-encoded, which imposes a numeric order that means nothing and probably understates how informative they are. TFT v2 fixes this.

\textbf{Attention and RQ1.} Attention peaks at \textbf{lag $-4$h}, with a cluster from $-2$h to $-5$h, which independently corroborates the Phase~1 lag OLS peak at $+6$h (Section~\ref{sec:lag-ols}): two very different methods both put the dominant news-impact horizon several hours after publication. There is also a secondary cluster at \textbf{$-27$h to $-28$h}, a ``daily memory'' effect at the same hour on the previous trading day, which the lag OLS could never have seen because it stops at 12h.

\begin{figure}[h]
  \centering
  \includegraphics[width=0.85\textwidth]{tft_attention_weights}
  \caption{TFT v1 encoder attention by lag. Mean attention over the 48-hour encoder window, peaking at $-4$h. Two clusters are visible: a short-term one between $-2$h and $-5$h, and a secondary one around $-27$h to $-28$h.}
  \label{fig:tftv1-attention}
\end{figure}

\textbf{Directional asymmetry (RQ2).} Splitting validation samples by sentiment direction gives the mean predicted \texttt{log\_volume} in Table~\ref{tab:tftv1-asymmetry}.

\begin{table}[h]
  \centering
  \caption{TFT v1 mean predicted \texttt{log\_volume} by sentiment direction.}
  \label{tab:tftv1-asymmetry}
  \begin{tabular}{@{}lr@{}}
    \toprule
    Sentiment direction & Mean predicted \texttt{log\_volume} \\
    \midrule
    Bearish & 8.775 \\
    Bullish & 8.700 \\
    Neutral & 8.875 \\
    \bottomrule
  \end{tabular}
\end{table}

Bearish predicted volume comes out above bullish, which points the same way as Section~\ref{sec:lag-ols}, but a two-sample t-test on the difference is not significant (\textbf{$p = 0.56$}) because a few hundred hours per category is not enough to power it. So the lag OLS stays the canonical RQ2 evidence and the TFT only corroborates the direction. Neutral hours have the highest predicted volume of the three, which is an ``uncertainty drives trading'' story we flag as speculative and leave for the discussion.

\textbf{Limitations of TFT v1 that drive TFT v2.} TFT v1 validates the approach: LLM sentiment dominates, attention corroborates the lag OLS, and the asymmetry shows up in direction. But five limitations each motivate a Phase~2 change:

\begin{enumerate}
  \item \textbf{Underpowered asymmetry test} ($p = 0.56$), motivating a larger corpus.
  \item \textbf{Integer-encoded categoricals} impose a meaningless numeric order, motivating proper categorical encoding.
  \item \textbf{No entity awareness} (entities extracted but unused), motivating 71 entity flags.
  \item \textbf{Pre-war training data only}, motivating a corpus extended through May 2026.
  \item \textbf{The dominant feature is the unreliable composite sentiment}: the 0.39-correlation feature from Section~\ref{sec:calibration} carries 53\% of the model's weight, motivating the channel decomposition.
\end{enumerate}

The fifth one is the pivot. The model leans hardest on exactly the feature the calibration flagged as least trustworthy, so v2 replaces it with the orthogonal channels (Section~\ref{sec:channel-response}) and retrains (Section~\ref{sec:tft-v2}).

\subsection{Channel decomposition response}\label{sec:channel-response}

The calibration (Section~\ref{sec:calibration}) showed the composite \texttt{sentiment\_score} was unreliable ($r = 0.39$, 31\% sign disagreement) on exactly the high-magnitude geopolitical events that matter most, and TFT v1 was putting 53\% of its weight on it. This subsection is the schema revision that fixes that, and the re-calibration that checks the fix.

\textbf{Design principle.} An oil-market article carries up to three economically distinct judgments: what it implies for \textbf{supply}, for \textbf{demand}, and for \textbf{geopolitical or operational risk}. That mirrors the standard structural separation of oil-price drivers into supply, aggregate-demand, and precautionary (risk) components \citep{kilian2009}. A single composite forces all three into one number, and when they point different ways it breaks. A supply shock is bearish for the event but bullish for price. A de-escalation is positive in valence but bearish for supply through sanctions relief. Schema v2 therefore splits the judgment up so each score carries one factual claim.

\textbf{Schema v2 changes.} The main addition is the three orthogonal channels, each on $[-1, +1]$: \texttt{supply\_impact} (positive = more oil to market), \texttt{demand\_impact} (positive = stronger consumption), and \texttt{risk\_premium} (positive = elevated risk). Schema v2 also promotes \texttt{usable} to a required field, makes \texttt{event\_type} a salience-ordered array, drops the redundant \texttt{price\_direction}, and keeps the composite \texttt{sentiment\_score}. Appendix~\ref{app:schema} has the full field-by-field v1-vs-v2 comparison.

\textbf{Re-calibration.} Re-scoring the same 30-article sample with Schema v2 improves every metric (Table~\ref{tab:recalibration}).

\begin{table}[h]
  \centering
  \caption{Inter-model re-calibration metrics, Schema v1 versus v2.}
  \label{tab:recalibration}
  \begin{tabular}{@{}lll@{}}
    \toprule
    Metric & v1 schema & v2 schema \\
    \midrule
    \texttt{usable} agreement & 26/30 (87\%) & 27/30 (90\%) \\
    \texttt{sentiment\_score} correlation & 0.39 & \textbf{0.88} \\
    \texttt{sentiment\_score} sign disagreements & 4/13 usable (31\%) & \textbf{1/14 usable (7\%)} \\
    \texttt{supply\_impact} correlation & --- & 0.94 \\
    \texttt{supply\_impact} MAD & --- & 0.11 \\
    \texttt{supply\_impact} sign disagreements & --- & 0 \\
    \texttt{demand\_impact} correlation & --- & 0.96 \\
    \texttt{demand\_impact} MAD & --- & 0.05 \\
    \texttt{demand\_impact} sign disagreements & --- & 0 \\
    \texttt{risk\_premium} correlation & --- & 0.82 \\
    \texttt{risk\_premium} MAD & --- & 0.13 \\
    \texttt{risk\_premium} sign disagreements & --- & 1 \\
    \bottomrule
  \end{tabular}
\end{table}

The headline numbers: \texttt{sentiment\_score} correlation goes from \textbf{0.39 to 0.88} and its sign disagreements from \textbf{31\% to 7\%}, and every new channel correlates at 0.82 or better (\texttt{supply\_impact} 0.94, \texttt{demand\_impact} 0.96, \texttt{risk\_premium} 0.82) with near-zero sign disagreements. The surprising part is that the composite improved even though we never touched it. Asking the model to decompose the article first and only then produce a composite seems to discipline everything it outputs. The channels are also genuinely orthogonal (within-model pairwise $\lvert r\rvert < 0.5$), so they are capturing distinct economic dimensions rather than the same judgment under three names.

\textbf{Residual weakness.} On 3 of the 14 usable calibration articles, Haiku zeroed all three channels where GPT gave small non-zero values, on articles of modest relevance. The matching \texttt{sentiment\_score} and \texttt{magnitude} were subdued too, so these do not dominate anything downstream, but this conservatism is the same pattern the \texttt{usable\_strict} filter is built on (Section~\ref{sec:filter-comparison}). With the channels validated, the full Schema v2 batch of 19{,}619 articles produces the feature set for TFT v2 (Section~\ref{sec:tft-v2}).

\subsection{Temporal Fusion Transformer v2}\label{sec:tft-v2}

TFT v1 showed the deep-learning approach works. TFT v2 brings together everything Phase~2 invested in (the channel decomposition of Section~\ref{sec:channel-response}, the LLM filter of Section~\ref{sec:filter-comparison}, the expanded corpus, and entity normalization) and is the model we evaluate the research questions against. A controlled ablation guided the design (Appendix~\ref{app:ablation}), and the configuration reported here is the one that won on validation performance.

\subsubsection{Success criteria}\label{sec:tft-v2-criteria}

The three criteria reflect what Phase~2 is for. The Phase~1 OLS gives the primary evidence for RQ1 (the $+6$h peak) and RQ2 (bearish $>$ bullish), and TFT v2's job is to corroborate that through an independent method while showing the channels and entities as visible drivers.

\begin{enumerate}
  \item \textbf{Channels and entities are visible drivers.} The VSN importance ranking should place channel features and entity flags in the top ranks, with a non-degenerate, interpretable attention pattern.
  \item \textbf{Prediction substantially beats persistence} in the 1-to-12h window, confirming the Phase~2 features carry predictive signal.
  \item \textbf{The multi-horizon error curve is consistent with Phase~1's lag structure}, with the largest persistence reduction in the $+6$h range identified by the lag OLS.
\end{enumerate}

Section~\ref{sec:tft-v2-eval} has the evaluation.

\subsubsection{Model design}\label{sec:tft-v2-design}

TFT v2 uses a hidden size of 32, four attention heads, dropout 0.15, and a hidden-continuous size of 16, for 298{,}329 parameters. It trains on the \texttt{usable\_strict=1} filter (Section~\ref{sec:filter-comparison}), so on 10{,}514 articles, with the 1{,}161 channel-neutral ones excluded. The input features are:

\begin{itemize}
  \item \textbf{Channels (Phase 2 contribution):} \texttt{supply\_impact}, \texttt{demand\_impact}, \texttt{risk\_premium}, each on $[-1, +1]$.
  \item \textbf{Composite news:} \texttt{sentiment\_score}, \texttt{magnitude}, \texttt{certainty}, \texttt{n\_articles}.
  \item \textbf{Categoricals:} \texttt{event\_type\_primary} (8 categories including the synthetic \texttt{no\_news}) and \texttt{time\_horizon} (4), as proper embeddings (dimension 5 and 3), taking the dominant article's value (highest \texttt{magnitude}, tie-broken on \texttt{article\_id}).
  \item \textbf{Entity flags (Phase 2 contribution):} 71 binary multi-hot columns (Section~\ref{sec:llm-extraction}).
  \item \textbf{Market context:} \texttt{log\_volume}, \texttt{price\_range}, \texttt{log\_return}, \texttt{amihud}, \texttt{dxy}, \texttt{vix}.
  \item \textbf{Calendar:} \texttt{hour}, \texttt{day\_of\_week}, \texttt{month}, \texttt{is\_us\_session}, \texttt{is\_wednesday}.
\end{itemize}

Continuous features are hour-averaged, entity flags take the maximum, and categoricals take the highest-magnitude article. Hours with no news get the \texttt{no\_news} category and zeros. A few boundary nulls (DXY/VIX holiday gaps, first-hour \texttt{log\_return}/\texttt{amihud}) are forward-filled or set to zero, which affects under 0.1\% of samples.

The model predicts three targets jointly (\texttt{log\_volume}, \texttt{amihud}, \texttt{price\_range}) at four horizons (\textbf{1, 3, 6, 12h}), and the $+6$h horizon is there to match the lag OLS peak (Section~\ref{sec:lag-ols}). The temporal split is \textbf{60/20/20} on the 11{,}232 hourly observations, with 48h buffers between partitions: train 2024-05-13 to 2025-06-11 (6{,}728 samples), validation 2025-06-13 to 2025-12-27 (2{,}216), test 2025-12-29 to 2026-05-13 (2{,}159). The war onset (28 February 2026) lands inside the test window, splitting it into 994 pre-war and 1{,}165 war hours, so the test set measures how a pre-war-trained model generalizes to an unseen war regime. We report metrics on the full test set, the pre-war slice, and the war slice throughout.

Training uses Adam (learning rate 1e-3, halved after three epochs without validation improvement). The objective is a quantile loss applied to each of the three targets independently and then summed, so the model learns a spread of quantiles (a prediction interval) per target rather than a single point estimate. Each target is normalized separately, with normalization statistics pooled within the asset group, which here is just WTI, so the three targets sit on comparable scales before their losses combine. Gradients are clipped at 0.1, training stops early once validation loss flattens (patience 10, minimum delta 1e-4), and the run is seeded (42) with deterministic algorithms on a Colab T4 GPU. Best validation loss was \textbf{0.427} at epoch 10.

\subsubsection{Predictive performance}\label{sec:predictive-performance}

We evaluate TFT v2 on the held-out test set for all three targets at all four horizons, against a persistence baseline (the current value carried forward, the standard financial-forecasting reference), reporting the median quantile (q50) on four slices: validation, full test, the pre-war test portion (994 hours before the 28 February 2026 onset), and the war portion (1{,}165 hours from onset to the end of the data).

\begin{table}[h]
  \centering
  \caption{\texttt{log\_volume}: TFT v2 versus persistence MAE.}
  \label{tab:mae-logvolume}
  \begin{tabular}{@{}rrrrrr@{}}
    \toprule
    Horizon & Persistence MAE & TFT v2 MAE & Reduction & Pre-war MAE & War MAE \\
    \midrule
    1h & 1.076 & 0.585 & 46\% & 0.536 & 0.628 \\
    3h & 1.452 & 0.577 & 60\% & 0.537 & 0.611 \\
    6h & 1.820 & 0.602 & 67\% & 0.551 & 0.646 \\
    12h & 2.174 & 0.631 & 71\% & 0.568 & 0.685 \\
    \bottomrule
  \end{tabular}
\end{table}

The error reduction over persistence climbs steadily with horizon, from 46\% at 1h to 71\% at 12h, which happens because persistence gets worse at longer horizons as volume autocorrelation decays. The $+6$h and $+12$h peak sits in the same range the lag OLS identified (Section~\ref{sec:lag-ols}). Error runs 17 to 21\% higher on the war slice than pre-war at every horizon, which is the price of extrapolating from pre-war-only training, but it is small enough that the model still has real predictive value on a regime it never saw.

\begin{table}[h]
  \centering
  \caption{\texttt{amihud}: TFT v2 versus persistence MAE.}
  \label{tab:mae-amihud}
  \begin{tabular}{@{}rrrrrr@{}}
    \toprule
    Horizon & Persistence MAE & TFT v2 MAE & Reduction & Pre-war MAE & War MAE \\
    \midrule
    1h & 0.0004 & 0.00023 & 43\% & 0.00010 & 0.00035 \\
    3h & 0.0004 & 0.00023 & 43\% & 0.00010 & 0.00034 \\
    6h & 0.0004 & 0.00022 & 45\% & 0.00009 & 0.00034 \\
    12h & 0.0004 & 0.00022 & 45\% & 0.00009 & 0.00033 \\
    \bottomrule
  \end{tabular}
\end{table}

Amihud gives a steady 43 to 45\% reduction across horizons with little variation, though part of that flatness is an artifact of amihud's compressed 1e-4 scale. War-slice persistence MAE is three to four times the pre-war value, which reflects illiquidity spikes in the war regime, but the model's relative reduction over persistence holds up on both slices.

\begin{table}[h]
  \centering
  \caption{\texttt{price\_range}: TFT v2 versus persistence MAE.}
  \label{tab:mae-pricerange}
  \begin{tabular}{@{}rrrrrr@{}}
    \toprule
    Horizon & Persistence MAE & TFT v2 MAE & Reduction & Pre-war MAE & War MAE \\
    \midrule
    1h & 0.495 & 0.718 & $-45$\% & 0.180 & 1.200 \\
    3h & 0.578 & 0.719 & $-24$\% & 0.180 & 1.200 \\
    6h & 0.630 & 0.721 & $-14$\% & 0.180 & 1.198 \\
    12h & 0.701 & 0.724 & $-3$\% & 0.180 & 1.201 \\
    \bottomrule
  \end{tabular}
\end{table}

This is where the model loses. It underperforms persistence on the full test set, by $-3$ to $-45$\%, and the failure is specific to a regime rather than general. On the pre-war slice it matches persistence at an MAE of about 0.180. On the war slice it degrades to about 1.200, which is worse than war-slice persistence at 0.764 to 1.059. Having only ever seen the moderate-volatility pre-war regime, the model reverts toward the historical mean when asked to predict the war regime, while persistence at least tracks the elevated current state. This is a clean regime-extrapolation failure, and it makes sense that price\_range is where it shows up: it measures intraday volatility directly, so it is far more regime-sensitive than the aggregate targets. Chapter~\ref{chap:discussion} takes up the implications.

\textbf{Across targets.} The model gets substantial gains on \texttt{log\_volume} (46 to 71\%) and \texttt{amihud} (43 to 45\%) and fails on \texttt{price\_range} in the war regime, which fits the train/test regime mismatch: the aggregate liquidity measures transfer across regimes and \texttt{price\_range} does not. So the rest of Section~\ref{sec:tft-v2} concentrates on \texttt{log\_volume}, which is both the target most relevant to RQ1 and RQ2 and the one the model handles best. The \texttt{price\_range} failure comes back in Section~\ref{sec:tft-v2-lag} and in Chapter~\ref{chap:discussion}.

\subsubsection{Feature contributions}\label{sec:feature-importance}

The Variable Selection Network (VSN) gives each input feature a learned weight for how much the model relies on it, and aggregating those weights over validation samples ranks the features. Table~\ref{tab:feature-importance} has the top ten for \texttt{log\_volume} prediction.

\begin{table}[h]
  \centering
  \caption{TFT v2 top ten features by VSN importance (\texttt{log\_volume}).}
  \label{tab:feature-importance}
  \begin{tabular}{@{}rlrl@{}}
    \toprule
    Rank & Feature & Importance & Type \\
    \midrule
    1 & \texttt{vix} & 0.188 & macro covariate \\
    2 & \texttt{supply\_impact} & 0.121 & channel (news) \\
    3 & \texttt{ent\_oman} & 0.113 & entity flag (news) \\
    4 & \texttt{demand\_impact} & 0.055 & channel (news) \\
    5 & \texttt{is\_wednesday} & 0.022 & calendar \\
    6 & \texttt{ent\_japan} & 0.017 & entity flag (news) \\
    7 & \texttt{ent\_eu} & 0.016 & entity flag (news) \\
    8 & \texttt{ent\_iran} & 0.015 & entity flag (news) \\
    9 & \texttt{ent\_china} & 0.014 & entity flag (news) \\
    10 & \texttt{ent\_algeria} & 0.014 & entity flag (news) \\
    \bottomrule
  \end{tabular}
\end{table}

Three patterns come out of this. First, VIX ranks first at 18.8\%, which fits its role as a market-wide volatility proxy: \texttt{log\_volume} tracks broad equity-market activity through stress periods, and risk expectations shape how much volume responds to oil news. Second, the channel decomposition puts two features in the top five, \texttt{supply\_impact} at 12.1\% and \texttt{demand\_impact} at 5.5\%, which together are 17.6\% of importance, roughly matching VIX on its own. The composite \texttt{sentiment\_score} does not make the top ten at all, which suggests that once the model has the decomposed channels, the aggregate sentiment has little left to add. Third, six of the top ten are entity flags, which validates the entity normalization of Section~\ref{sec:llm-extraction}, and they are all economically readable: \texttt{ent\_oman} is a producer next to the Strait of Hormuz, \texttt{ent\_iran} is the actor at the center of the conflict, \texttt{ent\_japan}, \texttt{ent\_eu}, and \texttt{ent\_china} are major importers, and \texttt{ent\_algeria} is a mid-size OPEC+ producer.

The two feature-engineering choices behind this result, proper categorical encoding and the entity flags, were checked in a controlled three-variant ablation. Its key finding is that the channels' predictive role only appears under proper categorical encoding, where \texttt{demand\_impact} becomes the top feature, and not under the integer encoding of the simplest variant, where the composite sentiment score dominates instead. The full three-variant comparison is in Appendix~\ref{app:ablation}, to keep the main text on the canonical model.

The VSN tells us a feature's weight, not the direction or size of its effect, so this reading is qualitative: VIX informs volatility expectations, the channels capture fundamental supply and demand news, and the entity flags carry source-specific risk signals. Sections~\ref{sec:tft-v2-lag} and~\ref{sec:asymmetry} look at how the model turns these features into predictions, through its attention pattern and its per-horizon error structure.

\textbf{Stability of the ranking.} The exact ordering in Table~\ref{tab:feature-importance} does not survive retraining, and it is worth being direct about that. An independent run, retrained after adding two candidate covariates and a minor data-cleaning fix, reshuffled the ranking substantially: \texttt{risk\_premium} and \texttt{magnitude} moved into the top two while \texttt{supply\_impact} and \texttt{ent\_oman} dropped out of it. What is stable across both runs is the character of the leading features rather than their order. Risk and volatility proxies (VIX, \texttt{risk\_premium}) and news-salience measures lead in both, while the composite \texttt{sentiment\_score} and the article count sit near the bottom in both. So these importance results should be read at the level of feature type, meaning risk and salience over price direction (Section~\ref{sec:asymmetry} and Chapter~\ref{chap:discussion}), and not as a precise per-feature ranking.

\subsubsection{Lag structure analysis: evidence for RQ1}\label{sec:tft-v2-lag}

RQ1 asks at what lag news events most strongly affect liquidity. Phase~1's lag OLS (Section~\ref{sec:lag-ols}) put the bearish-news peak on \texttt{log\_volume} at $+6$h with a secondary trace at $+12$h. We test TFT v2 against that with two diagnostics: the per-horizon error curve and the attention pattern.

\textbf{Per-horizon error curve.} If news impact really concentrates at a particular horizon, then a multi-horizon model should be most accurate there once we normalize for how hard each horizon is to begin with, which is what MAE-reduction-over-persistence does. On \texttt{log\_volume} (Table~\ref{tab:mae-logvolume}) that curve grows steadily from 46\% at $+1$h to 71\% at $+12$h, putting the strongest persistence-relative improvement in the $+6$h to $+12$h window, the same range the lag OLS found. The exact peak differs a little, $+12$h for TFT v2 against $+6$h for the lag OLS. The two methods are answering structurally different questions: lag OLS isolates the direct effect of one sentiment value at lag $-k$, while the TFT predicts forward from the full 48-hour window using all features at once. That is why their agreement on this window counts as corroboration through an independent methodology rather than the same finding said twice.

\textbf{Attention pattern.} The encoder attention distribution, aggregated over validation samples, gives a second and complementary view of the lag structure.

\begin{figure}[h]
  \centering
  \includegraphics[width=0.85\textwidth]{attention_by_sentiment_tftv2.2-exp2}
  \caption{TFT v2 mean encoder attention across the 48-hour window, disaggregated by sentiment direction. Overall attention peaks at $-1$h; bearish-sentiment hours peak at $-6$h while bullish-sentiment hours peak at $-1$h.}
  \label{fig:attention-sentiment}
\end{figure}

Figure~\ref{fig:attention-sentiment} shows mean encoder attention across the 48-hour window. Overall attention peaks at lag $-1$h (weight 0.0315) and falls off smoothly through $-2$h to $-5$h (0.0308, 0.0304, 0.0301, 0.0300), a graded recency effect with nothing piling up at the boundary. Split by sentiment direction, though, the pattern comes apart: bearish-sentiment hours peak at $-6$h (0.0321) while bullish-sentiment hours peak at $-1$h (0.0347). The bearish peak at $-6$h lines up exactly with the $+6$h lag OLS peak. That says the model has learned bearish news needs a longer integration window than bullish news, and it confirms independently the same structure Phase~1 found in the average bearish impact.

\textbf{Synthesis on RQ1.} Both diagnostics back the lag OLS finding from methodologically different angles: the per-horizon error curve peaks in the $+6$h to $+12$h range, and the bearish-sentiment attention peaks precisely at $-6$h. Between them the $+6$h lag has substantive empirical support from two independent analyses.

\subsubsection{Directional asymmetry analysis: evidence for RQ2}\label{sec:asymmetry}

RQ2 asks whether bearish news produces a different liquidity response than comparable bullish news. Phase~1's lag OLS (Section~\ref{sec:lag-ols}) found a robust bearish $>$ bullish asymmetry in \texttt{log\_volume}, strongest at $+6$h, and it remains the primary answer to RQ2. What we test here is whether TFT v2's predictions add complementary evidence through an independent method.

\textbf{Test design.} At each horizon and slice we split samples by the sentiment score at the prediction hour into bearish (below $-0.1$), bullish (above $+0.1$), and neutral (within plus or minus 0.1), then test whether mean predicted \texttt{log\_volume} differs between the bearish and bullish subsets using Welch's t-test with unequal variance. Neutral hours are excluded because they capture low-news activity, dominate the sample counts, and dilute the contrast the test is built to find. That gives sixteen tests, four horizons times four slices (val, test\_full, test\_prewar, test\_war), in Table~\ref{tab:asymmetry-tests}.

\begin{table}[h]
  \centering
  \small
  \caption{Directional asymmetry tests: mean predicted \texttt{log\_volume}, bearish versus bullish.}
  \label{tab:asymmetry-tests}
  \begin{tabular}{@{}rlrrrrrrl@{}}
    \toprule
    Horizon & Slice & Bear.\ $n$ & Bull.\ $n$ & Bear.\ mean & Bull.\ mean & Diff.\ & $p$ & Sig. \\
    \midrule
    1h & val & 320 & 271 & 8.345 & 8.281 & $+0.064$ & 0.578 & No \\
    1h & test\_full & 384 & 826 & 8.928 & 8.857 & $+0.071$ & 0.316 & No \\
    1h & test\_prewar & 180 & 340 & 8.853 & 8.783 & $+0.070$ & 0.543 & No \\
    1h & test\_war & 204 & 486 & 9.001 & 8.901 & $+0.100$ & 0.253 & No \\
    3h & val & 320 & 271 & 7.938 & 8.131 & $-0.193$ & 0.270 & No \\
    3h & test\_full & 384 & 826 & 8.732 & 8.596 & $+0.136$ & 0.247 & No \\
    3h & test\_prewar & 180 & 340 & 8.706 & 8.365 & $+0.341$ & 0.058 & No \\
    3h & test\_war & 204 & 486 & 8.756 & 8.729 & $+0.027$ & 0.869 & No \\
    6h & val & 320 & 271 & 8.186 & 8.369 & $-0.183$ & 0.230 & No \\
    6h & test\_full & 384 & 826 & 8.633 & 8.529 & $+0.104$ & 0.400 & No \\
    6h & test\_prewar & 180 & 340 & 8.470 & 8.380 & $+0.089$ & 0.636 & No \\
    6h & test\_war & 204 & 486 & 8.788 & 8.615 & $+0.174$ & 0.285 & No \\
    12h & val & 320 & 271 & 8.277 & 8.371 & $-0.094$ & 0.561 & No \\
    12h & test\_full & 384 & 826 & 8.220 & 8.464 & $-0.244$ & 0.070 & No \\
    12h & test\_prewar & 180 & 340 & 8.274 & 8.333 & $-0.059$ & 0.761 & No \\
    12h & test\_war & 204 & 486 & 8.169 & 8.540 & $-0.371$ & 0.053 & No \\
    \bottomrule
  \end{tabular}
\end{table}

\textbf{Results.} None of the sixteen tests reaches significance at $p < 0.05$. The closest is $+12$h on test\_war (bearish 8.169 vs bullish 8.540, difference $-0.371$, $p = 0.053$). Two others get near without crossing: $+3$h test\_prewar (bearish $>$ bullish, $+0.341$, $p = 0.058$) and $+12$h test\_full (bearish $<$ bullish, $-0.244$, $p = 0.070$).

\textbf{Interpretation.} This test and Phase~1's lag OLS estimate different quantities, and that distinction is what determines how to read the null. The OLS compares the marginal coefficients on bearish and bullish sentiment, so the sensitivity of volume to each direction with the other held fixed (Section~\ref{sec:lag-ols}). The test here compares the mean predicted volume of bearish hours against bullish hours, which is an unconditional difference of group means with nothing held fixed at all. A larger marginal sensitivity to bearish news, which is what Phase~1 found, can sit perfectly well alongside near-equal average volume across bearish and bullish hours, for instance if bullish hours happen to cluster in higher-baseline-volume periods like risk-on sessions. So the absence of a group-mean asymmetry here does not contradict Phase~1. It answers a coarser, confounded question.

The point-prediction differences are small and non-significant at every horizon, with no consistent sign across slices. This is not a power problem: the bearish and bullish groups each hold hundreds of samples, and no test becomes significant under alternative sentiment thresholds swept from 0.0 to 0.5 once we correct for multiple comparisons. What it actually mirrors is the target. In the real test data, the future volume of bearish and bullish hours is statistically indistinguishable at the 1h to 6h horizons, and the only significant ground-truth difference, at $+12$h, has bullish above bearish, a sign the model reproduces with a predicted difference of $-0.244$. So the model is tracking a group-level near-symmetry that is genuinely there in its target period, not washing out a signal that exists. The one directional structure the TFT does carry is temporal rather than in magnitude: the attention pattern (Section~\ref{sec:tft-v2-lag}) goes to $-6$h for bearish sentiment and $-1$h for bullish, matching Phase~1's $+6$h peak.

\textbf{Synthesis on RQ2.} Phase~1's lag OLS stays the primary and appropriate evidence for RQ2, because it measures the directional sensitivity the research question is actually about. TFT v2 shows no magnitude asymmetry in its point predictions, and that is the expected outcome rather than a failed replication: the group-mean comparison is a different and confounded estimand, and the model's target is itself near-symmetric across sentiment direction over the test period. What the model genuinely adds is temporal, the attention asymmetry of Section~\ref{sec:tft-v2-lag}, where bearish is integrated over a longer window ($-6$h) than bullish ($-1$h). That corroborates the $+6$h lag structure of RQ1 without leaning on the magnitude test at all.

\subsubsection{Success criteria evaluation}\label{sec:tft-v2-eval}

We now check TFT v2 against the three success criteria declared in Section~\ref{sec:tft-v2-criteria}, using the evidence from Sections~\ref{sec:predictive-performance} through~\ref{sec:asymmetry}.

\begin{table}[h]
  \centering
  \small
  \caption{TFT v2 success-criteria evaluation.}
  \label{tab:success-criteria}
  \begin{tabular}{@{}p{4.5cm}lp{7.5cm}@{}}
    \toprule
    Criterion & Status & Evidence \\
    \midrule
    1. Channels and entities as visible drivers & Pass & \texttt{supply\_impact} and \texttt{demand\_impact} appear in the top four features (12.1\% and 5.5\%); six entity flags in the top ten; attention layer non-degenerate with peak at $-1$h; bearish/bullish attention divergence ($-6$h vs $-1$h) matches Phase~1's identified lag \\
    2. Substantial persistence-relative improvement & Pass & \texttt{log\_volume} MAE reduction of 46--71\% across the 1--12h window and \texttt{amihud} 43--45\% (\texttt{price\_range} is the exception on the war regime, Section~\ref{sec:predictive-performance}) \\
    3. Multi-horizon curve consistent with Phase 1 & Pass & Persistence reduction peaks at $+12$h (71\%) with $+6$h close behind (67\%); attention pattern independently confirms the $+6$h finding via the bearish-sentiment attention peak \\
    \bottomrule
  \end{tabular}
\end{table}

TFT v2 passes all three. The Phase~2 features, the channels and entities, are economically interpretable and drive prediction; the model substantially beats persistence on the aggregate liquidity targets; and its multi-horizon error curve and attention independently corroborate Phase~1's $+6$h lag structure. None of that makes Phase~2 the source of evidence for the research questions. Phase~1's lag OLS is still the primary statistical evidence for RQ1 (the $+6$h peak) and RQ2 (the bearish $>$ bullish sensitivity), while TFT v2 contributes evidence that is complementary and methodologically distinct: interpretable feature importance, and a bearish-sentiment attention peak at $-6$h confirming the same lag structure. The two phases are complementary, not in competition. The configuration evaluated here is what the rest of the thesis means by TFT v2.

\subsection{Methodological comparison: Phase 1 versus Phase 2}\label{sec:phase-comparison}

The two phases come at the same research questions through structurally different methods. Phase~1 (Section~\ref{sec:phase1}) uses regression, contemporaneous and lag OLS, on FinBERT sentiment over a regex-filtered corpus. Phase~2 (Section~\ref{sec:phase2}) uses a deep-learning forecaster, TFT v2, on LLM-derived channel-decomposed features over an LLM-filtered corpus. Neither replaces the other. They answer complementary questions about the same data.

\subsubsection{What Phase 2 changed methodologically}

Phase~2 changed the analytical stack at five levels (corpus, feature extraction, tone versus price impact, entity normalization, and analytical framework), and each change answers a specific Phase~1 limitation.

\textbf{Corpus expansion.} Phase~1's 13{,}690-article regex-filtered corpus (\texttt{body\_valid}) only covers the pre-war period to February 2026. Phase~2's LLM-filtered corpus (\texttt{usable\_strict=1}, 10{,}514 articles) runs through May 2026, so it includes the 28 February 2026 war onset and what followed, and the \texttt{usable\_strict} filter (Section~\ref{sec:filter-comparison}) drops 1{,}161 topical-but-channel-neutral articles. The Phase~2 corpus is smaller in count but more complete in time, because it covers a regime change Phase~1's data simply does not contain.

\textbf{Feature extraction.} Phase~1's FinBERT gives one three-class sentiment label per article (bearish, neutral, bullish), averaged hourly. Phase~2's LLM extracts a structured schema per article: a composite sentiment score, three channels (\texttt{supply\_impact}, \texttt{demand\_impact}, \texttt{risk\_premium}), magnitude, certainty, event type, time horizon, and entities. The channels are what let the model separate supply, demand, and risk-premium news that FinBERT flattens onto one axis.

\textbf{Tone versus price impact.} The two sentiment signals do not just differ in resolution, they measure different things, and the difference turns out to be diagnostic. FinBERT classifies the surface financial tone of an article. The Haiku \texttt{sentiment\_score} is prompted for the net directional impact on the WTI price, with the instruction to ``assess the actual content, not the headline tone''. Across the 6{,}300 articles scored by both methods the two agree only moderately (Pearson 0.54), and the disagreement is concentrated in one place: geopolitical news, where mean FinBERT tone is negative ($-0.116$) while mean Haiku price impact is bullish ($+0.137$). On demand, macro, and supply articles the two agree in sign. The mechanism is visible in the headlines they split on, for example ``Russian Sanctions Drive Oil Prices Higher'' (FinBERT $-0.95$, Haiku $+0.85$) and ``Oil prices soar after Israel launches strikes on Iran capital'' (FinBERT $-0.96$, Haiku $+0.80$). FinBERT anchors on the alarming vocabulary, sanctions and strikes, and labels the article negative even when the headline says outright that prices rose, because it has no asset grounding and no way to reason from a supply threat to a higher price. The LLM, prompted to think about price impact for WTI specifically, reads the same articles as bullish. This is the FinBERT limitation that motivated the Phase~2 extraction, in its most concrete form, and it is confined to exactly the geopolitical news carrying the highest risk premium and the highest trading volume (Section~\ref{sec:feature-importance}). It also explains how a sentiment-sign asymmetry test can come out differently across the two phases (Section~\ref{sec:asymmetry}): the high-volume geopolitical news lands in opposite sentiment buckets depending on which definition you use.

\textbf{Entity normalization.} Phase~1 has no entity features at all. Phase~2 maps entity mentions onto 71 canonical flags (Section~\ref{sec:llm-extraction} and Appendix~\ref{app:entities}), giving each hour a binary vector of which actors (US, Iran, Saudi Arabia, OPEC+, and so on) the dominant article named. No Phase~1 method had access to that signal.

\textbf{Analytical framework.} Phase~1's regression isolates the effect of specific variables at specific lags but produces no forecast. Phase~2's TFT predicts three targets at four horizons jointly from the previous 48 hours, but cannot isolate individual variable effects the way OLS does. This is not a move from a weaker method to a stronger one. It is a move from one that isolates effects to one that integrates them, and each one covers the other's blind spot.

\subsubsection{Why cross-phase numbers are not comparable}

TFT v1 and TFT v2 loss values cannot be compared directly, because nearly everything underneath them is different: the loss function (v1 \texttt{QuantileLoss} on one target, v2 \texttt{MultiLoss} on three, so roughly three times the per-target scale), the validation period (v1 roughly 2{,}160 pre-war hours from an 80/20 split, v2 roughly 2{,}216 transitional hours from the 60/20/20 split), the horizons (v1 one hour, v2 four jointly), the feature count (v1 about 12, v2 over 90), and the upstream corpus and extraction (FinBERT on the regex corpus versus LLM channels on the \texttt{usable\_strict} corpus). Putting v1's val\_loss of 0.204 next to v2's 0.427 would therefore mislead rather than inform.

The fair comparison is persistence-relative reduction on a shared target and horizon (\texttt{log\_volume} at 1h), where both models beat persistence substantially and neither fails. Phase~2's contribution is not better forecasts, it is a richer analytical vocabulary: where v1 dumps its importance onto one scalar sentiment score, v2 spreads it across decomposed channels, entity flags, and macro covariates. So the progression from Phase~1 to Phase~2 is real but qualitative and analytical, not a numerical win. The thesis' primary contribution is the Phase~2 methodology, not a claim that it forecasts better.

\section{Summary}\label{sec:experiments-summary}

Chapter~\ref{chap:experiments} has reported the empirical results of both phases.

Phase~1 (Section~\ref{sec:phase1}) established the initial evidence for the research questions, through regression on a pre-war corpus of 13{,}690 articles with FinBERT sentiment features. The headline bias experiment (Section~\ref{sec:headline-bias}) showed that title-only and title-plus-body FinBERT inputs produce substantially different sentiment classifications, which motivated the move to richer feature extraction in Phase~2. The lag OLS (Section~\ref{sec:lag-ols}), including the contemporaneous $k=0$ case, confirmed a positive news-volume association and put the peak of news impact on liquidity at the $+6$ hour lag, with a statistically significant bearish $>$ bullish asymmetry. An exploratory vector autoregression (Section~\ref{sec:var}) was abandoned once the sparsity of the hourly sentiment signal proved incompatible with its assumptions, which reinforced the move to the deep-learning models of Phase~2.

Phase~2 (Section~\ref{sec:phase2}) extended the analysis with LLM-based feature extraction on a corpus covering both the pre-war and war regimes. The Schema v2 extraction produced channel-decomposed features (Section~\ref{sec:llm-extraction}) separating supply-side, demand-side, and risk-premium news impacts, alongside 71 canonical entity flags from a normalized entity list. Comparing the regex-based Phase~1 filter against the LLM-based Phase~2 filter (Section~\ref{sec:filter-comparison}) identified specific error modes in each and motivated the stricter \texttt{usable\_strict=1} variant used for TFT v2 training. TFT v1 (Section~\ref{sec:tft-v1}) validated the deep-learning approach on a single target and horizon, beating persistence substantially and showing a two-cluster attention pattern (short-term at $-4$h, daily memory at $-27$/$-28$h). TFT v2 (Section~\ref{sec:tft-v2}) integrated the full set of Phase~2 investments: channel decomposition, entity flags, proper categorical encoding, multi-horizon prediction, and multi-target learning. The v2 canonical configuration cuts prediction error over persistence by 46 to 71 percent on \texttt{log\_volume} across horizons, peaking at $+12$h, consistent with the $+6$ to $+12$ hour range Phase~1's lag analysis identified. The Variable Selection Network ranks VIX and \texttt{supply\_impact} as the top two features and puts six entity flags in the top ten. The directional asymmetry test, sixteen bearish-versus-bullish comparisons of predicted volume, returned no significant group-level difference. That reflects a different and coarser estimand than Phase~1's marginal-coefficient asymmetry rather than a failed replication, and the TFT's complementary evidence for RQ2 is temporal: its attention goes to $-6$h for bearish sentiment against $-1$h for bullish.

The methodological comparison of the two phases (Section~\ref{sec:phase-comparison}) argued that comparing v1 and v2 numbers directly is not appropriate, given the differences in validation sets, loss functions, prediction horizons, feature counts, and news corpus. Both models beat persistence substantially in their own settings, and Phase~2's contribution is analytical richness rather than measurably better raw accuracy. That section also established that FinBERT's tone-based sentiment and the LLM's price-impact sentiment measure different constructs, agreeing overall but diverging in sign on the high-volume geopolitical news.

The success criteria set for TFT v2 before training were evaluated in Section~\ref{sec:tft-v2-eval}, and all three pass. The evaluation refines what the Phase~1 findings mean through a different methodological lens, and it exposes the \texttt{price\_range} regime-extrapolation limitation that feeds into the discussion. Chapter~\ref{chap:discussion} develops the implications of these findings along with the limitations and directions for future work. Chapter~\ref{chap:conclusion} concludes.
